\section{Conclusion}

We opened with a concrete problem: multi-dimensional trustworthiness evaluation can miss compound failures invisible to independent assessment. A medical diagnosis LLM might achieve 92\% on truthfulness and 89\% on fairness when evaluated separately, yet fail on both dimensions for the same minority-group patients---a coupling pattern current benchmarks cannot detect. Our methodology framework makes these hidden vulnerabilities measurable.

Through phi coefficient analysis demonstrated on synthetic benchmark emulation across 5 trustworthiness dimensions and 3 simulated model profiles, we establish that coupling is detectable when present (phi 0.33--0.40, $p < 1\times10^{-13}$) but measurably sparse---limited to 2 dominant dimension pairs, not pervasive across all 10 possible pairs. Truthfulness-robustness coupling (phi 0.36--0.40) appears consistently across GPT-4, Claude-3, and Llama-3 profiles, while fairness-safety coupling (phi 0.33--0.40) shows comparable strength. No model exhibits coupling in 3 or more pairs after correction for multiple comparisons, establishing sparsity as a detectable property rather than a measurement artifact.

This sparsity detection capability persists when we control for instance difficulty. Partial correlation analysis yields partial phi values of 0.36--0.56 with effect size retention of 86--161\%---in five of six cases, controlling for difficulty strengthens rather than weakens coupling in synthetic data. Hard instances do not drive these patterns in our validation; the methodology isolates shared underlying vulnerabilities from confounds. If validated on real benchmarks, coupling would reflect architectural properties: calibration failures linking truthfulness and robustness, value alignment training linking fairness and safety.

Models show distinct coupling profiles in synthetic data---GPT-4 shows a cognitive coherence chain (truthfulness$\rightarrow$robustness$\rightarrow$safety), Claude-3 shows a value alignment cluster (fairness$\rightarrow$safety$\rightarrow$privacy), Llama-3 shows minimal coupling---though statistical confirmation requires larger samples than our proof-of-concept provides. These observational fingerprints suggest methodology could detect different vulnerability architectures reflecting distinct training approaches, not universal coupling patterns.

If validated on real benchmarks, the implications are immediate. Multi-dimensional evaluation frameworks should report coupling statistics alongside per-dimension scores. Benchmark designers should ensure coverage of detected coupling pairs for compound risk assessment. Model selection teams could exploit coupling patterns: prioritize models with strong coupling in required dimensions when deployment scenarios demand multiple trustworthiness guarantees simultaneously.

Looking forward, real-world validation with production benchmarks (MultiTrust, TrustLLM) and production models remains the critical next step. Our synthetic proof-of-concept validates the measurement approach but leaves real-world coupling magnitudes unknown. Mechanistic interpretation---why these two pairs couple---offers a path to principled intervention: if calibration failures drive truthfulness-robustness coupling in real models, targeted calibration improvements should weaken that link. A coupling profile database covering 10+ models would enable deployment-oriented model selection tools, turning coupling awareness from research methodology into operational practice. Temporal tracking across model versions could reveal whether updates alter vulnerability architectures, providing a new lens for monitoring model evolution.

Trustworthiness dimensions are neither universally coupled---sharing a single architectural bottleneck that fails broadly---nor universally independent---exhibiting zero correlation across all pairs. Our methodology demonstrates they can exhibit sparse coupling, a detectable property requiring targeted mitigation strategies rather than universal solutions. Where coupling exists, it reveals shared vulnerabilities. Where coupling is absent, dimensions fail independently. Both patterns matter. Both are measurable via this framework; real-world magnitudes remain unknown pending production benchmark access.
